\section{Sujet et place dans le service}
\subsection{Sujet}
\subsubsection{Sujet défini avant mon arrivée}
~\\
Avant le début de mon stage, il était prévu que je travaille sur un projet innovant visant à développer une preuve de concept pour extraire des informations numériques sur les émissions de carbone à partir de rapports de durabilité en pdf. 
Pour ce faire, nous devions utiliser des technique de Natural Language Processing en python.\\
Le soutien aux tâches en cours afin de connaître toutes les fonctions clés de l'équipe faisait également partie intégrante du stage. 
Cela aurait pu inclure la documentation des données, la liaison et la comparaison des ensembles de données pour la gestion de la qualité des données et la participation aux groupes de travail internes et externes.


\subsubsection{Sujet réel}
~\\
Mes tâches en réalité étaient plus variées et sur divers projets :
\begin{itemize}
    \item Crawling de résultats de recherche Google pour récupérer les rapports de durabilité des entreprises, et extraction d'informations des PDF
    \item Optimisation de code Python pour la génération et comparaison de gros jeux de données
    \item Rapprochement de jeux de données provenant de différentes sources (analyses)
    \item Génération automatisée d'un rapport à partir de résultats stockés dans un fichier csv
    \item Projet de recherche NLP pour trouver quels datasets et méthodologies sont utilisé(e)s dans quels papiers de recherche
\end{itemize}


\pagebreak

\subsection{Fonction occupée dans le service}
~\\
Ma fonction dans le département IT de l'association était double. 
J'étais avant tout développeuse web et devait donc me charger de l'implémentation de diverses fonctionnalités du site web, en front comme en back, le tout en ReactJS et Symfony.
\\
Mais j'ai également occupé la fonction de coordinatrice de niveau 3, 
c'est-à-dire que je gérais un sous-projet avec l'aide d'un collègue, en l'occurence le projet Répertoire Vert. Mes missions étaient les suivantes :
\begin{itemize}
    \item Accueillir les nouveaux arrivants
    \item Avoir une vision globale des tâches à accomplir et évaluer leur difficulté
    \item Distribuer les tâches à chaque nouveau sprint en tenant compte des objectifs de formation et des capacités de chacun
    \item Faire un planning pour respecter les deadlines
    \item Vérifier le travail de l'équipe
    \item Mettre le site en production deux fois par semaine
    \item Présenter l'avancée au client
\end{itemize}


% % Use H to place the figure HERE
% \begin{figure}[H]
%     \centering
%     \includegraphics[width=\textwidth]{latexsup.png}
%     \caption{Quod Erat Demonstrandum}
% \end{figure}

% % to prevent the picture to go at the end of this part
% \FloatBarrier

% \clearpage

% \subsection{Une sous-section}

% Une liste :
% \begin{itemize} % liste sans ordre
%     \item Niveau 1 - \href{https://fr.wikipedia.org/wiki/USB}{USB}
%     \begin{itemize}
%         \item Niveau 2 - \href{https://fr.wikipedia.org/wiki/Ethernet}{Ethernet}
%         \begin{itemize}
%             \item Un élément de niveau 3 - IP
%             \begin{enumerate} % Liste ordonnée
%                 \item Un élément de niveau 4 - \href{https://en.wikipedia.org/wiki/TCP}{TCP}
%                 \item Un second élément de niveau 4 - \href{https://en.wikipedia.org/wiki/UDP}{UDP}
%             \end{enumerate}
%         \end{itemize}
%         \item Retour au niveau deux - \href{https://fr.wikipedia.org/wiki/Spanning_Tree_Protocol}{STP}
%     \end{itemize}
%     \item[(NomChoisi)] Un autre élément de niveau 1 - \href{https://fr.wikipedia.org/wiki/Carrier_Sense_Multiple_Access_with_Collision_Avoidance}{CSMA/CA}
% \end{itemize}

% Une liste sur deux colonnes :
% \begin{itemize}[twocol]
%     \item Je suis sur la première colonne !
%     \item Moi aussi.
%     \item Et moi ?
%     \item Je ne sais pas moi.
%     \item Toi tu es de l'autre côté, sur la première colonne !
%     \item Moi, je suis avec toi sur la deuxième colonne.
% \end{itemize}


% \subsection{Une autre sous-section}
% Pour savoir comment faire de longs tableaux, vous pouvez vous référer à l'Annexe \ref{longtab}
% \subsubsection{Une sous-sous-section}
% Un excellent professeur proclama un jour:
% \begin{center}
% Il fait trop chaud pour faire du réseau.
% \end{center}

% A l'\textbf{extrême gauche} on a:
% \begin{flushleft}
%     Coucou comment ça va ?
% \end{flushleft}

% Tandis qu'à l'\underline{extrême droite} on a le \href{https://rassemblementnational.fr/}{\footnote{Rassemblement National}{RN}} et aussi cette mise en forme:

% \begin{flushright}
%     Vous ne trouvez pas que petit, on a tous voulu changer la société avant que ce soit elle qui nous change ?
% \end{flushright}

% \subsubsection{Une autre sous-sous-section}
% \paragraph{Un paragraphe}
% Une citation c'est bien, mais bien citer c'est mieux :

% \begin{quoting}
%     Mais, vous savez, moi je ne crois pas qu’il y ait de bonne ou de mauvaise situation. Moi, si je devais résumer ma vie aujourd’hui avec vous, je dirais que c’est d’abord des rencontres, des gens qui m’ont tendu la main, peut-être à un moment où je ne pouvais pas, où j’étais seul chez moi. Et c’est assez curieux de se dire que les hasards, les rencontres forgent une destinée… Parce que quand on a le goût de la chose, quand on a le goût de la chose bien faite, le beau geste, parfois on ne trouve pas l’interlocuteur en face, je dirais, le miroir qui vous aide à avancer. Alors ce n’est pas mon cas, comme je le disais là, puisque moi au contraire, j’ai pu ; et je dis merci à la vie, je lui dis merci, je chante la vie, je danse la vie… Je ne suis qu’amour ! Et finalement, quand beaucoup de gens aujourd’hui me disent : « Mais comment fais-tu pour avoir cette humanité ? » Eh bien je leur réponds très simplement, je leur dis que c’est ce goût de l’amour, ce goût donc qui m’a poussé aujourd’hui à entreprendre une construction mécanique, mais demain, qui sait, peut-être simplement à me mettre au service de la communauté, à faire le don, le don de soi…
%     \begin{flushright}
%         --- Otis, Astérix Mission Cléopatre
%     \end{flushright}
% \end{quoting}

% Si vous appréciez la façon "Markdown" de présenter les citations, je vous propose la même chose ici :
% \begin{quoted}
%     Ceci est une citation comme usuellement vue sur \texttt{Notion} ou en \texttt{Markdown}.
% \end{quoted}

% \subparagraph{Un sous-paragraphe}

% \begin{dialogue}
%     \speak{Un Allemand} \direct{s'esclaffe} Tous les allemands ne sont pas Nazis !
%     \speak{Hubert Bonisseur de La Bath} Oui, je connais cette théorie
% \end{dialogue}

% \section*{Une section non numérotée}
% \addcontentsline{toc}{section}{Une section non numérotée} % Si on veut qu'elle apparaisse dans le sommaire.
% On peut créer une mise en forme attirant l'attention sur un point important à expliquer :

% \begin{callout}{Contrôle de flux $\neq$ contrôle de congestion}
%     \begin{itemize}
%         \item Le \textbf{contrôle de flux} signifie essentiellement que TCP s'assure qu'un expéditeur ne submerge pas un destinataire en envoyant des paquets plus vite qu'il ne peut les consommer. Il concerne le nœud final.
%         \item Le \textbf{contrôle de congestion} vise à empêcher un nœud de submerger le réseau (c'est-à-dire les liens entre deux nœuds).
%     \end{itemize}
% \end{callout}

% Ou plus sobrement :

% \begin{myboxedtext}
%     Avoir un joli rapport $\Rightarrow$ $+50$ points de charisme.
% \end{myboxedtext}